\chapter{Généralités sur les fonctions}
\thispagestyle{fancy}

\vspace*{-0.4cm}
\section{Définitions et notations}

\vspace*{-0.4cm}
\subsection{Définition}

\begin{bexemple}
	On considère la fonction \(f\) qui exprime l'aire d'un rectangle de dimensions 3 et \(x\).\\
	Une expression littérale de \(f\) est donc : \(f(x)=3 x\).
\end{bexemple}
\begin{bdefi}
	\phantom{$\dfrac{1}{1}$}\pointille[\textwidth]
	
	\phantom{$\dfrac{1}{1}$}\pointille[\textwidth]
\end{bdefi}

\vspace*{-0.4cm}
\subsection{Image et antécédent}

\begin{bexemple}
	Dire que : \(f(2)=5\) signifie que : 
	
	\includegraphics[scale=0.25]{Cours_img01}
	
	On dit que :
	\begin{itemEx}
		\item l'image de 2 par la fonction \(f\) est 5.
		\item un antécédent de 5 par \(f\) est 2.
	\end{itemEx}
\end{bexemple}

\begin{brmqs}
	\begin{itemRem}
		\item Un nombre possède une unique image.
		\item Cependant, un nombre peut posséder plusieurs antécédents.
	\end{itemRem}
\end{brmqs}

\begin{bmethodet}{Déterminer l'image d'une fonction par calcul}
	\grasInMet{Enoncé:}
	
	Soit la fonction \(g\) définie par \(g(x)=x^{2}-2\).\\
	Calculer l'image de 6 par la fonction \(g\).
	
	\grasInMet{Réponse:}
	
	\phantom{$\dfrac{1}{1}$}\pointille[\textwidth]
	\phantom{$\dfrac{1}{1}$}\pointille[\textwidth]
	\phantom{$\dfrac{1}{1}$}\pointille[\textwidth]
\end{bmethodet}


\begin{bmethodet}{Déterminer un antécédent par calcul}
	\grasInMet{Enoncé:}
	
	Soit la fonction \(f\) définie par \(f(x)=2 x-3\).\\
	Déterminer un antécédent de $-5$ par la fonction \(f\).
	
	\grasInMet{Réponse:}
	
	\phantom{$\dfrac{1}{1}$}\pointille[\textwidth]
	\phantom{$\dfrac{1}{1}$}\pointille[\textwidth]
	\phantom{$\dfrac{1}{1}$}\pointille[\textwidth]
	\phantom{$\dfrac{1}{1}$}\pointille[\textwidth]
	\phantom{$\dfrac{1}{1}$}\pointille[\textwidth]
	\phantom{$\dfrac{1}{1}$}\pointille[\textwidth]
\end{bmethodet}

\vspace*{-1.2cm}
\section{Représentation graphique}

\vspace*{-0.6cm}
\begin{bmethodet}{Représenter graphiquement une fonction}
	\grasInMet{Enoncé:}
	
	Soit la fonction \(f\) définie par \(f(x)=5 x-x^{2}\).\\
	On donne un tableau de valeurs de la fonction \(f\) :
	\begin{center}
		\begin{tabular}[t]{|l|l|l|l|l|l|l|l|l|}
			\hline
			\(x\) & 1 & 1,5 & 2 & 2,5 & 3 & 3,5 & 4 & 4,5 \\
			\hline
			\(f(x)\) & 4 & 5,25 & 6 & 6,25 & 6 & 5,25 & 4 & 2,25 \\
			\hline
		\end{tabular}
	\end{center}
	Tracer, dans un repère, la courbe représentative de la fonction \(f\).
	
	\grasInMet{Réponse:}
	
	\begin{minipage}[c]{0.48\linewidth}
		\phantom{$\dfrac{1}{1}$}\pointille[\textwidth]
		\phantom{$\dfrac{1}{1}$}\pointille[\textwidth]
		\phantom{$\dfrac{1}{1}$}\pointille[\textwidth]
		\phantom{$\dfrac{1}{1}$}\pointille[\textwidth]
		\phantom{$\dfrac{1}{1}$}\pointille[\textwidth]
		\phantom{$\dfrac{1}{1}$}\pointille[\textwidth]
		\phantom{$\dfrac{1}{1}$}\pointille[\textwidth]
		\phantom{$\dfrac{1}{1}$}\pointille[\textwidth]
		\phantom{$\dfrac{1}{1}$}\pointille[\textwidth]
	\end{minipage} \hfill
	\begin{minipage}[c]{0.48\linewidth}
		\includegraphics[scale=0.25]{Cours_img02}
	\end{minipage}
	
	
	
\end{bmethodet}

\vspace*{-1.2cm}
\section{Résolution graphique d'équations et d'inéquations}

\vspace*{-0.6cm}
\begin{bmethodet}{Résoudre graphiquement une équation}
	\grasInMet{Enoncé:}
	
	\begin{minipage}[c]{0.48\linewidth}
		On a représenté la courbe de la fonction \(f\) définie par \(f(x)=5 x-x^{2}\).\\
		Résoudre graphiquement l'équation \(5 x-x^{2}=4\).\\
	\end{minipage} \hfill
	\begin{minipage}[c]{0.48\linewidth}
		\includegraphics[scale=0.2]{Cours_img03}
	\end{minipage}
	
	\newpage
	
	\grasInMet{Réponse:}
	
	\begin{minipage}[c]{0.48\linewidth}
		\phantom{$\dfrac{1}{1}$}\pointille[\textwidth]
		\phantom{$\dfrac{1}{1}$}\pointille[\textwidth]
		\phantom{$\dfrac{1}{1}$}\pointille[\textwidth]
		\phantom{$\dfrac{1}{1}$}\pointille[\textwidth]
		\phantom{$\dfrac{1}{1}$}\pointille[\textwidth]
		\phantom{$\dfrac{1}{1}$}\pointille[\textwidth]
		\phantom{$\dfrac{1}{1}$}\pointille[\textwidth]
	\end{minipage} \hfill
	\begin{minipage}[c]{0.48\linewidth}
		\includegraphics[scale=0.25]{Cours_img04}
	\end{minipage}
	
	
\end{bmethodet}

\begin{brmqs}
	\begin{itemRem}
		\item Par lecture graphique, les solutions obtenues sont approchées.
		\item L'équation \(f(x)=7\), par exemple, ne semble pas avoir de solution car la courbe représentée ne possède pas de point d'ordonnée 7.
		\item Graphiquement, on ne peut pas être certain que les solutions qui apparaissent sont les seules. Il pourrait y en avoir d'autres au-delà des limites de la représentation graphique tracée.
	\end{itemRem}
\end{brmqs}

\begin{bmethodet}{Résoudre graphiquement une inéquation}
	\grasInMet{Enoncé:}
	
	Dans la méthode précédente, on a représenté la courbe de la fonction \(f\) définie par \(f(x)=5 x-x^{2}\).\\
	Résoudre graphiquement l'inéquation \(5 x-x^{2}>4\).
	
	\grasInMet{Réponse:}
	
	\begin{minipage}[c]{0.48\linewidth}
		\phantom{$\dfrac{1}{1}$}\pointille[\textwidth]
		\phantom{$\dfrac{1}{1}$}\pointille[\textwidth]
		\phantom{$\dfrac{1}{1}$}\pointille[\textwidth]
		\phantom{$\dfrac{1}{1}$}\pointille[\textwidth]
		\phantom{$\dfrac{1}{1}$}\pointille[\textwidth]
		\phantom{$\dfrac{1}{1}$}\pointille[\textwidth]
	\end{minipage} \hfill
	\begin{minipage}[c]{0.48\linewidth}
		\includegraphics[scale=0.23]{Cours_img05}
	\end{minipage}
	
	
	
\end{bmethodet}

\vspace*{-1.0cm}
\section{Variations d'une fonction}

\vspace*{-0.5cm}
\subsection{Taux de variation}

\begin{bdefi}
	\phantom{$\dfrac{1}{1}$}\pointille[\textwidth]
	\phantom{$\dfrac{1}{1}$}\pointille[\textwidth]
\end{bdefi}

\begin{bprop}
	\phantom{$\dfrac{1}{1}$}\pointille[\textwidth]
	\phantom{$\dfrac{1}{1}$}\pointille[\textwidth]
\end{bprop}

\vspace*{-0.8cm}
\begin{bmethodet}{Déterminer un taux de variation d'une fonction}
	\grasInMet{Enoncé:}
	
	Soit \(f\) la fonction définie sur \(\mathbb{R}\) par : \(f(x)=2 x^{2}+1\).
	\vspace*{-0.2cm}
	\begin{enumMet}
		\item Déterminer le taux de variation entre 1 et 3.
		\item Interpréter géométriquement ce taux de variation.
	\end{enumMet}
	\vspace*{-0.2cm}
	\grasInMet{Réponse:}
	
	\vspace*{-0.2cm}
	\begin{minipage}[c]{0.68\linewidth}
		\begin{enumMet}
			\item \phantom{$\dfrac{1}{1}$}\pointille[\textwidth]
			\phantom{$\dfrac{1}{1}$}\pointille[\textwidth]
			\phantom{$\dfrac{1}{1}$}\pointille[\textwidth]
			\phantom{$\dfrac{1}{1}$}\pointille[\textwidth]
			\item \phantom{$\dfrac{1}{1}$}\pointille[\textwidth]
			\phantom{$\dfrac{1}{1}$}\pointille[\textwidth]
			\phantom{$\dfrac{1}{1}$}\pointille[\textwidth]
		\end{enumMet}
	\end{minipage} \hfill
	\begin{minipage}[c]{0.28\linewidth}
		\includegraphics[scale=0.14]{Cours_img06}
	\end{minipage}
	
	
\end{bmethodet}
\vspace*{-1.2cm}
\subsection{Fonctions monotones}
\vspace*{-0.3cm}
\begin{bdefi}
	On dit qu'une fonction \(f\) est monotone sur un intervalle I, si \(f\) est :
	\begin{itemDef}
		\item \phantom{$\dfrac{1}{1}$}\pointille[\textwidth-1cm]
		\item \phantom{$\dfrac{1}{1}$}\pointille[\textwidth-1cm]
		\item \phantom{$\dfrac{1}{1}$}\pointille[\textwidth-1cm]
	\end{itemDef}
\end{bdefi}
%\newpage

\vspace*{-0.6cm}
\begin{bprops}
	\begin{itemProp}
		\item Si le taux de variation d’une fonction f entre deux nombres quelconques d’un intervalle I est
		\phantom{$\dfrac{1}{1}$}\pointille[\textwidth-1cm]
		\item \phantom{$\dfrac{1}{1}$}\pointille[\textwidth-1cm]
		\item \phantom{$\dfrac{1}{1}$}\pointille[\textwidth-1cm]
	\end{itemProp}
\end{bprops}

\vspace*{-0.3cm}
\begin{bmethodet}{Étudier les variations d'une fonction à l'aide du taux de variation}
	\vspace*{-0.2cm}
	\grasInMet{Enoncé:}
	
	Soit \(f\) la fonction définie sur \(\mathbb{R}\) par : \(f(x)=5 x-3\).\\
	Démontrer que \(f\) est strictement croissante sur \(\mathbb{R}\).
	
	\grasInMet{Réponse:}
	
	\phantom{$\dfrac{1}{1}$}\pointille[\textwidth]
	\phantom{$\dfrac{1}{1}$}\pointille[\textwidth]
	\phantom{$\dfrac{1}{1}$}\pointille[\textwidth]
	\phantom{$\dfrac{1}{1}$}\pointille[\textwidth]
\end{bmethodet}
